\documentclass[10pt]{article}

\usepackage[utf8]{inputenc}

\usepackage[T1]{fontenc}

\usepackage{amsmath}

\usepackage{amsfonts}

\usepackage{amssymb}

\usepackage[version=4]{mhchem}

\usepackage{stmaryrd}

\usepackage{bbold}



\begin{document}

произвольным нётеровым кольцам. Понятие П. к. позволяет в аксиоматич. форме учесть важнейшие из этих свойств.



Аксиомы П. к. $A$.



А1. Кольцо $A$ универсально цепное. (Кольцо $A$ наз. ц е п н ы м, если для любых двух его простых идеалов $\mathfrak{p} \neq \mathfrak{p}^{\prime}$ длины любых неуплотняемых цепочек простых идеалов $\mathfrak{p}=\mathfrak{p}_{0} \subset \mathfrak{p}_{1} \subset \ldots \subset \mathfrak{p}_{n}=\mathfrak{p}^{\prime} \quad$ совпадают. Кольцо $A$ наз. у н и в е р с а л ь н о це п ны м, если депным является любое кольцо многочленов $A\left[T_{1}, \ldots, T_{k}\right]$.)



А2. Формальные слои кольца $A$ являются геометрически регулярными, т. е. для любого простого идеала $\mathfrak{p} \subset A$ и гомоморфизма из $A$ в поле $K$ кольцо $\widehat{A}_{\mathfrak{p}} \bigotimes_{A} K$ регулярно. Здесь $\hat{A}_{\mathfrak{p}}$ - пополнение локального кольца $A_{\mathfrak{p}}$.



A3. Для любой целостной конечной $A$-алгебры $B$ найдется ненулевой элемент $b \in B$ такой, что кольцо частных $B\left[b^{-1}\right]$ регулярно.



П. к. обладают следующими свойствами.



\begin{enumerate}

  \item Для П. к. $A$ множество регулярных (соответственно нормальных) точек схемы $\operatorname{Spec} A$ открыто.



  \item Если превосходное локальное кольцо $A$ приведенное (соответственно нормальное, равноразмерное), то таким же будет пополнение $\hat{A}$.



  \item Целое замыкание П. к. $A$ в конечном расширении поля частных кольца $A$ является конечной $A$-алгеброй.



  \item Если кольцо $A$ превосходное, то любая $A$-алгебра ковечного типа - также П. к.



\end{enumerate}



Два важнейших примера П. к.- полные локальные кольца (или аналитич. кольца) и дедекиндовы кольца с полем частных нулевой характеристики. Таким образом, класс П. к. достаточно птирок и, в частности, содержит все алгебры конечного типа над полем или над кольцом целых чисел $\mathbb{Z}$.



Превосходность кольца $A$ тесно связана с возможностью разрешения особенностей схемы Spec $A$ (см. [1], [2]).



Jum.: [1] G rothendieck A. [D i e u donné J.], Eléments de géométrie algébrique, pt. 2, P., 1965; [2] X и p oн а к а X., “Математика", 1965, т. 9 , № 6, с. 3-70.



ПРЕДБАЗА - семейство $\gamma$ открытых $\begin{aligned} & \text { п. И. Данилов. } \\ & \text { подожеств }\end{aligned}$ топологич. пространства $X$ такое, что совокупность всех множеств, являющихся пересечением конечного числа әлементов $\gamma$, образует базу $X$.



ПРЕДВАРЕННАЯ ФОРМУЛА - ф. И. Войщеховский. исчисления предикатов (УИП), имеющая вид



\[

Q_{1} x_{1} Q_{2} x_{2} \ldots Q_{n} x_{n} \Psi

\]



где $Q_{i}$ обозначает квантор всеобщности $\forall$ или квантор существования $\exists$, переменные $x_{i}, x_{j}$ различны при $i \neq j$ и $\Psi$ - формула, не содержацая кванторов. П. Ф. наз. также п р е д в а $\mathrm{p}$ өны м и но р м а ль ь ы и и фо рмами или прене кстым и формами. Для всякой формулы $\varphi$ языка УИП можно найти П. Ф., логически әквивалентную в классич. исчислении предикатов формуле $\varphi$. Процесс нахождения П. Ф. основан на следуюгих әквивалентностях, выводимых в классич. УИП:



\[

\begin{gathered}

(\forall x \varphi(x) \supset \Psi) \equiv \exists x^{\prime}\left(\varphi\left(x^{\prime}\right) \supset \Psi\right), \\

\exists x \varphi(x) \supset \Psi \equiv \forall x^{\prime}\left(\varphi\left(x^{\prime}\right) \supset \Psi\right), \\

\Psi \supset \forall x \varphi(x) \equiv \forall x^{\prime}\left(\Psi \supset \varphi\left(x^{\prime}\right)\right), \\

\Psi \supset \exists x \varphi(x) \equiv \exists x^{\prime}\left(\Psi \supset \varphi\left(x^{\prime}\right)\right), \\

7 \forall x \varphi \equiv \exists x\urcorner \varphi, \neg \exists x \varphi \equiv \forall x\rceil \varphi, \\

Q y \forall x \varphi \equiv \forall x \varphi, Q y \exists x \varphi \equiv \exists x \varphi,

\end{gathered}

\]



где $x^{\prime}$ - любая переменная, не входящая свободно в формулы $\varphi(x), \Psi$, и $\varphi\left(x^{\prime}\right)$ получается из $\varphi(x)$ заменой всех свободных вхождений $x$ на $x^{\prime}$; перемевная $y$ не входит свободно в $\forall x \varphi, \exists x \varphi$. Чтобы применять приведенные әквивалентности, надо предварительно выразить логич. связки через $\supset$ и 7 ; затем, применяя әти әквивалентности, постепенно передвигать все кванторы влево. Получающаяся в результате П. Ф. наз. п р е дв а p e в н о й ф о р м о й данной формулы.



Лит.: [1] М е н д е л ь с о н Э., Введение в математическую логику, пер. с англ., М., 1971 . В. Н. Гришин



ПРЕДГИЛЬБЕРТОВО ПРОСТРАНСТВО - векторное пространство $E$ над полем комплексных или действительных чисел, снабженное скалярным произведением $E \times E \rightarrow \mathbb{C}, x \times y \rightarrow(x, y)$, удовлетворяющим следующим условиям:



\begin{enumerate}

  \item $(x+y, z)=(x, z)+(y, z),(\lambda x, y)=\lambda(x, y),(y, x)=$ $=\overline{(x, y)}, x, y, z \in E, \lambda \in \mathbb{C}(\mathbb{R})$;



  \item $(x, x) \geqslant 0$ для $x \in E$.



\end{enumerate}



На П. п. $E$ определена полунорма $\|x\|=(x, x)^{\frac{1}{2}}$. Пополнение $E$ по этой полунорме является гильбертовылм пространством.



ПРЕДЕЛ - одно из основных понятий математики, означающее, что какая-то переменная, зависящая от другой переменной, при определенном изменении последней, неограниченно приближается к нек-рому постоянному значению. Основным при определении П. является понятие близости рассматриваемых объектов: только после его введения П. приобретает точный смысл. С П. связаны основные понятия математич. анализа: непрерывность, нроизводная, дифференциал, интеграл. Одним из простейших случаев П. является П. последовательности.



П редел последовательности. Пусть $X$ - топологич. пространство. Последовательность его точек $x_{n}, n=1,2, \ldots$, наз. сходящейся к точке $x_{0} \in X$ или, что то же самое, точка $x_{0}$ наз. п р е де л о м д а нной пос ле д о в а те льност и, если для любой окрөстности $U$ точки $x_{0}$ существует такое натуральное $N$, что для всех $n>N$ выполняется включение $x_{n} \in U$; при этом пишут



\[

\lim _{n \rightarrow \infty} x_{n}=x_{0}

\]



В случае, когда $X-$ хаусдорфово пространство, П. последовательности $x_{n} \in X, n=1,2, \ldots$, если он существует, единствен. Для метрич. пространства $X$ точка $x_{0}$ является П. последовательности $\left\{x_{n}\right\}$ тогда и только тогда, когда для любого $\varepsilon>0$ существует такое натуральное $N$, что для всех номеров $n>N$ выполняется неравенство $\rho\left(x_{n}, x_{0}\right)<\varepsilon$, где $\rho\left(x_{n}, x_{0}\right)-$ расстояние между точками $x_{n}$ и $x_{0}$. Если последовательность точек метрич. пространства сходится, то она ограничена. Последовательность точек полного метрич. пространства является сходящейся в том и только в том случае, когда она фундаментальная. В частности, это верно для числовых последовательностей, для к-рых исторически впервые возникло понятие П. последовательности. Для числовых последовательностей справедливы формулы



\[

\lim _{n \rightarrow \infty} c=c, \lim _{n \rightarrow \infty} c x_{n}=c \lim _{n \rightarrow \infty} x_{n},

\]



$c$ - произвольное фиксированное число,



\[

\begin{aligned}

\lim _{n \rightarrow \infty}\left(x_{n}+y_{n}\right) & =\lim _{n \rightarrow \infty} x_{n}+\lim _{n \rightarrow \infty} y_{n}, \\

\lim _{n \rightarrow \infty} x_{n} y_{n} & =\lim _{n \rightarrow \infty} x_{n} \lim _{n \rightarrow \infty} y_{n},

\end{aligned}

\]



а если $\lim _{n \rightarrow \infty} y_{n} \neq 0$, то



\[

\lim _{n \rightarrow \infty} \frac{x_{n}}{y_{n}}=\frac{\lim _{n \rightarrow \infty} x_{n}}{\lim _{n \rightarrow \infty} y_{n}} .

\]



Эти свойства числовых последовательностей переносятся на П. последовательностей более общих структур,





\end{document}